\documentclass[UTF8,fontset=none,zihao=-4,a4paper]{ctexart}
\usepackage{geometry}
\geometry{margin=2.6cm}
\usepackage{xeCJK}
\usepackage{fontspec}
\usepackage{tabularx}
\usepackage{booktabs}
\usepackage{longtable}
\usepackage{enumitem}
\usepackage{xcolor}
\usepackage[most]{tcolorbox}
\usepackage{titlesec}
\usepackage{fancyhdr}
\usepackage{hyperref}
\hypersetup{colorlinks=true,linkcolor=black,urlcolor=black,
  pdftitle={在风景里读历史——论黄裳的文史游记},pdfauthor={ajzhanghk}}

% --- CJK fonts (Linux: AR PL + WenQuanYi) ---
\setCJKmainfont{AR PL SungtiL GB}
\setCJKsansfont{WenQuanYi Zen Hei}
\setCJKmonofont{WenQuanYi Zen Hei Mono}
\setCJKfamilyfont{zhkai}{AR PL KaitiM GB}
\newcommand{\kai}{\CJKfamily{zhkai}}
\setmonofont{DejaVu Sans Mono}[Scale=0.9]

% --- route symbols missing from the Latin font through WenQuanYi ---
\usepackage{newunicodechar}
\newfontfamily\symbolfont{WenQuanYi Zen Hei}
\newunicodechar{★}{{\symbolfont\char"2605}}
\newunicodechar{○}{{\symbolfont\char"25CB}}
\newunicodechar{●}{{\symbolfont\char"25CF}}
\newunicodechar{·}{{\symbolfont\char"00B7}}
\newunicodechar{↔}{{\symbolfont\char"2194}}
\newunicodechar{→}{{\symbolfont\char"2192}}
\newunicodechar{←}{{\symbolfont\char"2190}}
\newunicodechar{≤}{{\symbolfont\char"2264}}
\newunicodechar{≥}{{\symbolfont\char"2265}}
\newunicodechar{–}{{\symbolfont\char"2013}}
\newunicodechar{※}{{\symbolfont\char"203B}}

% --- rare CJK names (李{\fb 昪}/李{\fb 璟}) absent from the Sung body font.
%     newunicodechar cannot catch them (xeCJK claims CJK codepoints first),
%     so they are wrapped with \fb (WenQuanYi) by a post-pass in the builder. ---
\setCJKfamilyfont{zhfb}{WenQuanYi Zen Hei}
\newcommand{\fb}{\CJKfamily{zhfb}}

\linespread{1.32}
\setlength{\parindent}{2em}
\setlist{itemsep=2pt,topsep=3pt,parsep=0pt}

% headings in sans (黑体)
\titleformat*{\section}{\Large\sffamily\bfseries}
\titleformat*{\subsection}{\large\sffamily\bfseries}
\titleformat*{\subsubsection}{\normalsize\sffamily\bfseries}

\pagestyle{fancy}
\fancyhf{}
\fancyhead[L]{\small\sffamily 在风景里读历史 · 论黄裳的文史游记}
\fancyhead[R]{\small\thepage}
\renewcommand{\headrulewidth}{0.3pt}

\begin{document}
\sloppy



\begin{titlepage}
\centering
\vspace*{3.2cm}
{\sffamily\bfseries\fontsize{28}{34}\selectfont 在风景里读历史}\\[10pt]
{\kai\Large 论黄裳的文史游记}\\[2.6cm]
{\large ——历史品味、文章趣味与风土人情}\\[0.6cm]
{\large 一篇关于黄裳「地点触发的文史随笔」的文艺评论}\\[3cm]
{\large 作者：ajzhanghk}\\[0.4cm]
{\normalsize 2026 年 6 月}\\
\vfill
\begin{tcolorbox}[colback=gray!6,colframe=gray!55,boxrule=0.6pt,
  arc=2pt,left=10pt,right=10pt,top=8pt,bottom=8pt,width=0.92\linewidth]
\small
\textbf{版权与引用说明}\quad 本文系\textbf{原创文艺评论}。黄裳（容鼎昌，1919--2012）
作品仍在版权保护期内（存续至 2062 年）；本文\textbf{不收录其受保护的作品全文}，
引述以公有领域的古人成句与篇名为主，凡涉黄裳本人文字者只取短句并归于评论性引用。
文中史实凡一时未据原书坐实者，以「据其文／疑／约」提示，严格标注见本项目数据库
\texttt{data/huangshang\_essays.json}。
\end{tcolorbox}
\end{titlepage}


\tableofcontents

\clearpage


\textbf{作者：ajzhanghk}


\begin{quote}\small\itshape
\textbf{题注}　本文是为《黄裳行旅散文集（重编）》撰写的导论性文艺评论，亦可独立成篇。文中凡涉黄裳具体篇目，均依据本项目已整理的篇目数据库（\texttt{data/huangshang\_essays.json}，五十九篇）及其原创内容摘要；凡史实考订一时未能据原书坐实者，行文中以"据其文""疑""约"等语提示，其严格的"确定／待核／推测"标注见数据库。黄裳作品（1919—2012）仍在版权保护期内，本文引述以公有领域的古人成句与篇名为主，凡涉黄裳本人文字者只取短句并归于评论性引用，不作大段移录。
\end{quote}

\medskip\hrule\medskip

\subsection{引子：他先看见的不是风景}


读黄裳的游记，常有一种被推迟的满足。你以为他要写鸡鸣寺的春色，他却先告诉你豁蒙楼已被"防空司令部"征用，写稿不成，只好在大殿里喝茶；你以为他要写莫愁湖的水光，他却把胜棋楼下朱元璋与徐达对弈的旧事、把南朝民歌里那个叫"莫愁"的女子，一并请了出来。风景在他笔下从不单独存在——它总是某段历史的入口，某个人物的背景，某一部旧书的注脚。


这并非偶然的趣味偏向，而是一种自觉的写作方法。黄裳的游记本质上不是旅游散文，而是\textbf{"地点触发的文史随笔"}：眼前的风景与旧迹只是引线，真正展开的是掌故、史事、书籍、戏曲、人物，而落点几乎总在历史兴亡之感。明乎此，才能解释一个表面上的悖论：一个写了一辈子"行旅"的人，他的文章里"走"的成分其实很轻，"读"的成分极重。他到一处，与其说是去看,不如说是去\textbf{核对}——核对书上读到的、史里记下的、戏中演过的那个地方，与眼前这个地方,究竟还剩下多少重叠。


本文想要论证的是：这种"即地读史"的写法，背后有一双非常具体的眼睛。这双眼睛\textbf{痴迷于一个特定的历史时刻——明清易代}；而它之所以痴迷于此，又因为黄裳本人写下那些最深沉的篇章时，正身处二十世纪中国自己的"易代之际"。\textbf{他读十七世纪，是为了看懂他亲历的二十世纪。} 历史品味、文章趣味、风土人情这三种常被分开谈论的特色，一旦放回这条"以古照今"的主轴上，便不再是并列的优点，而成为一个完整的、有内在动力的文学世界。


\medskip\hrule\medskip

\subsection{一、历史的眼光：为什么总是明清之际}


通读黄裳行旅诸篇，一个分布上的事实会逐渐浮现：他笔下反复回访的历史时刻，高度集中于\textbf{晚明、南明与清初}这一段。这不是题材的随机散布，而是一种几近偏执的精神引力。


请看他写南京。《王介甫与金陵》追的是北宋——王安石"骑驴粗服往来钟山"的晚境，以及那首被认作绝唱的《桂枝香·金陵怀古》（"诸公寄调桂枝香者三十余家，惟王介甫为绝唱"）；《南唐二陵》追的是更早的十世纪，李{\fb 昪}、李{\fb 璟}的钦陵顺陵，背后立着那位只做了不到四十年江山、却以亡国之音名垂词史的李后主。这两篇证明黄裳的历史视野并不窄。然而一旦进入他真正用力的腹地，时针几乎都停在十七世纪：《石巢园》写阮大铖的宅园——那座由造园名家计成设计、民间却讥为"裤子裆"的园子，写园主如何在繁华精丽之后降清失节；《扫叶楼》写画家龚贤，写他穿着僧服、作"扫叶"之状，以示不事新朝；《咏怀堂》牵出"明末四公子"之一的冒辟疆与《影梅庵忆语》里的董小宛；《关于吴梅村》写那位写下《圆圆曲》、却以仕清为终身之痛的贰臣诗人；《绝代的散文家——张宗子》推崇的张岱，正是以《陶庵梦忆》在故国倾覆后追忆繁华的遗民;《余淡心与金陵》钩沉的余怀，则以《板桥杂记》为秦淮旧梦立了一块碑。


把这一串名字排在一起，主题便不言自明：\textbf{它们几乎全是"易代之际，士人何为"的案例。} 阮大铖与龚贤构成一组最尖锐的对照——同样面对清兵入关，一个降，一个隐；一个把才情用在媚附，一个把才情守成气节。吴梅村则是第三种：既未殉国、又不甘心，于是把后半生过成一场漫长的自悔。黄裳面对古迹，最在意的从来不是兴废本身那点苍凉，而是\textbf{兴废关头一个具体的人做出了什么选择}。他的怀古是道德性的：一座园林的美，会因园主的失节而变味；一座楼的名，会因画家的风骨而被记住。这是把山水当作一份份"气节的考卷"来读。


为什么是十七世纪？因为明清易代是中国历史上"气节"这一命题被推到极致的时刻——它逼着每一个读书人当众交出答案：降，还是不降；仕，还是隐；忘,还是记。而这恰恰是黄裳写作金陵诸篇时——1946 至 1949——他自己和他那一代人正被反复追问的问题。这一点，下文还要详述。此处先确认一个判断：\textbf{黄裳的历史品味，不是博物馆式的，而是审判式的。} 他面对废墟时常有的，如本项目编辑原则所概括的，"不是凭吊的伤感，而是考据式的好奇与兴亡之际的清醒"。伤感会让人软下来，清醒却让他始终保持着一种打量的、甚至略带冷峻的距离。


更见功力的是，这种道德审判从不靠声色俱厉的议论完成，而是埋在\textbf{考据}里。他不说阮大铖可耻，他只把那座精美园子的设计者、那个不堪的诨名、以及园主最终的去向并置在一起，让事实自己说话。这是史家的笔法，也是他历史品味中最克制、最耐读的部分。


\medskip\hrule\medskip

\subsection{二、怀古之外，还有见证：记者黄裳与"易代"的现场}


如果黄裳只会怀古，他至多是一位趣味高雅的旧派文人。使他成为现代散文中一个不可替代的存在的，是他历史眼光的\textbf{另一半：他不只读死去的历史，他还亲手记录正在发生的历史。}


这一半来自他的记者身份。黄裳长期供职于上海《文汇报》，而他笔力最盛的金陵书写，正出自 1946 至 1949 那几年作为记者在南京的亲历。于是我们看到一组与怀古诸篇截然不同、却同样动人的文字：《老虎桥边看"知堂"》——他以记者身份，经延龄巷、国府路，一路找到老虎桥模范监狱，申请获准，赶在最终宣判之前探看了狱中的周作人；《泽存书库》——他记录汪伪要员陈群所建、藏书四十万册善本四千余部的那座书库（书库之名取《礼记》"手泽存焉"之义），如何在战后被接收、善本如何随之流散；《司徒雷登夹着皮包走了以后》——他以美国大使的黯然离华为切口，写下国民政府垮台前夕的政局，那篇文字与毛泽东同年的《别了，司徒雷登》恰成一种历史的互文；《解放后看江南》——1949 年秋天，他同时写下人民迎接解放的欢欣与劫后满目的疮痍，把一个时代转折的两面都按在了纸上。


请注意这组文字与上一节那组的\textbf{时间关系}。当黄裳在 1946—49 年的南京，亲眼看着一个政权"劫收"、腐烂、崩解，看着一座首都易主，看着像周作人这样的文人为当年的"落水"接受审判——他所亲历的，正是一场二十世纪的"易代"。而他偏偏又是那样一个对十七世纪易代之际念兹在兹的人。两者在他身上短路、接通了。\textbf{老虎桥里那个等待宣判的周作人，与他笔下降清的阮大铖、自悔的吴梅村，是同一道考题在不同世纪的重出。} 黄裳未必每次都点破这层对应，但他选择去写谁、怎么写，已经泄露了这副贯通古今的眼光：他用读十七世纪养成的眼光，去看自己身处的二十世纪；又反过来，因为亲历了一场易代，他对三百年前那场易代里的种种选择，有了同代人般的体贴与苛刻。


这种"怀古"与"见证"的双重性，是黄裳区别于一般文史散文家的根本。纯粹的怀古容易流于赏玩，纯粹的见证容易沦为时评；而黄裳让两者互相牵制——他的怀古因为有现实的重量而不轻浮，他的见证因为有历史的纵深而不短视。本项目把《司徒雷登》一篇定为"可选"而非"必收"，理由正是它"时评成分较重"；这恰恰从反面说明：当历史纵深退场、只剩当下评议时，黄裳那独有的味道就会变淡。他最好的文章，永远站在古与今的接缝处。


\medskip\hrule\medskip

\subsection{三、文章的趣味：在周作人的闲与鲁迅的硬之间}


谈黄裳而不谈他的文章本身，是不公平的，因为他首先是一位\textbf{文体家}。他的趣味有一个清晰的谱系，又有一种罕见的综合。


一面，他承的是\textbf{晚明小品到周作人的"闲笔"一脉}。这条线讲究的是从容、是旁逸斜出的掌故、是"文抄公"式地把古书旧籍信手拈来。最能见出这一师承的是《昆明杂记》：他在昆明购得一部记载南明永历帝流亡滇缅的晚明野史，意识到自己正踏在三百年前同一片土地上，于是索性用鲁迅所谓"抄古书"的笔法，把历史的断片抄录、缀连、点染成文。这种以征引旧籍结构全篇的写法，是知堂的看家本领，黄裳学得地道。他写江南山水也常如此——《采石·当涂·青山》里，他让陆游《入蜀记》对采石矶的描绘先行,再请出初到当涂、写下《望天门山》的李白，最后落到范传正将李白迁葬青山的史实上。三种文献、三个时代，被一条江水串成一篇,这就是"抄古书"抄出的境界：文章的厚度，来自纸背叠着的纸。


但黄裳若只有这"闲"的一面，便不过是又一个知堂的影子。他另有\textbf{鲁迅式的硬}。这"硬"，一在他前述的道德锋芒——他终究是要审判的；二在他的\textbf{反讽}。《关于"美国兵"》一组里，《S.O.S.》写美军后勤勤务部（Services of Supply），单是那个借自国际呼救信号的标题，反讽已先到——他以近乎正经的笔调铺陈其补给程序之繁复、物资之浪费、官僚机器之荒诞，越正经越好笑。《关于"翻译官"》更是夫子自道的一篇妙文：他辨析"翻译官"这个称呼的来历，指出外事局的公文写"X 级译员"，远征军里写"Interpreting Officer"，而民间则把这等人物比作"姨太太"、骂作"二毛子"——几个称谓排开，战时中外权力关系中那点尴尬、依附与自嘲，尽在其中。这种把自己也放进显微镜下的冷幽默，是周作人很少有、而鲁迅常有的。


于是黄裳的文章趣味，可以概括为一次成功的\textbf{调和：以周作人的从容，运鲁迅的骨力}。前者给他余裕与博雅，后者给他锋芒与脊梁。他能在一篇之内，上一段还在闲闲地考订一座楼的来历（豁蒙楼本是张之洞为纪念其遇难门生杨锐而建），下一段就让陈后主"胭脂井"的亡国旧典沉沉压来——《重过鸡鸣寺》一篇，把陈后主的亡国、十七世纪余怀父子的《金陵览古》、与 1980 年自己的重游叠成三重时空的对话，雅而不腐，深而不滞。雅与俗、闲与硬、考据与抒情，在他笔下被调得恰到好处。这是几十年功夫养出来的分寸感，学不来，也仿不像。


\medskip\hrule\medskip

\subsection{四、风土与人情：从茶馆的众生到一味"美人肝"}


黄裳的历史眼光若只向上看王侯将相、士人气节，他的世界会显得太"高"。幸而他还有结结实实向下、向俗的一面——这就是他笔下的\textbf{风土人情}。而值得玩味的是，他写风土，用的仍是那双"读史"的眼睛：一道菜、一间茶馆，在他看来同样是文化史的标本。


《茶馆》是其中的范本。他不写茶之雅，而写茶馆这个空间如何成为战时后方社会的横切面：广元、剑阁的茶叶粗砺，成都的茶馆松弛，重庆的茶馆甚至兼作谈判与交易之所；一室之中，警察可以与挑夫同座，大学生把这里当自修室,生意人把这里当交易所。寥寥几笔，一个阶层杂处、百业纷纭的战时社会便活了过来。这种写法的要害在于：他始终\textbf{从下往上看}。在《为美国兵活着的人们》里，他的镜头对准的是依附美军经济谋生的中国底层——街头摊贩、洗衣女工、酒吧女招待、翻译助手——他们的生计与尊严都系于那支外来军队的存在。这是抗战散文里不多见的、带着社会学冷静的阶级观察，其中的批判不靠口号，靠的是把这些人一个个看清楚、写下来。


最能见出黄裳"以小见史"功夫的，是《"美人肝"》。这道听来香艳、实为鸭胰所制的南京名馔，与松鼠鱼、蛋烧卖、凤尾虾并称"京苏菜四大名菜"，出自清真老字号马祥兴。黄裳写它，不只写口腹之乐，更写一道菜里沉淀的地方饮食传统与城市记忆——饮食在他这里也是一种"旧迹"，一种可供考订、可供怀想的文化遗存。一个能从"美人肝"里吃出文化史的人，与一个能从废园里读出气节的人，是同一个人。\textbf{他对世界的基本态度是：万物皆可读，处处都有史。} 风土人情之于黄裳，不是游记里调剂口味的边角料，而是历史在民间、在日常、在烟火气里的另一种存留。


这一面也让他那贯通古今的眼光避免了书斋的悬空。他记得周作人，也记得马祥兴的厨子；他追问阮大铖的气节，也心疼为美国兵洗衣服的女人。\textbf{历史在他笔下，既是士大夫的考卷，也是贩夫走卒的生计。} 唯其如此，他的兴亡之感才不是文人的自我感动，而扎根在一个有体温、有市声、有饭菜香气的真实中国里。


\medskip\hrule\medskip

\subsection{五、三事一体：行旅、读史、访书的合流}


要真正理解黄裳的独特，还须看到他把三件常被分开的事，合成了一件：\textbf{行旅即读史，读史即访书，而访书又是另一种行旅。} 这三者在他身上是一个循环,缺一不可。


他是著名的藏书家与古籍版本学者，他的"书话"里藏着大量行踪——这正是本项目专设第七辑"书肆、旧书与纸上行旅"的依据。《西泠访书记》写他在杭州孤山一带书肆间访书：书肆散落于湖山园林之间，于是访书的脚步同时也是游山的脚步，翻检旧册的过程同时也是阅读地理的过程。《琉璃厂》写北京那条旧书古玩街——从清咸丰年间始创的来薰阁，到 1925 年陈济川接掌后旧书业的格局——访书在这里又成了对一座城市文化记忆的考古。对黄裳来说，\textbf{登临是读史，访书是考古，二者本是一事}：都是穿过眼前的实物，去触摸背后那层层叠叠的时间。


《泽存书库》是这一"合流"的极致样本。那是一篇关于书的文章，也是一篇关于政治的文章，还是一篇关于战争与流散的文章——汪伪要员的搜罗、四十万册藏书的命运、善本最终随败退的政权远走台湾，书史与政治史在此完全咬合,再难分开。在黄裳手里，一部书的递藏，就是一段政权的兴亡；访求一部旧籍的路线图，就是一幅历史的地形图。这是把"书旅"正式纳入"行旅"的有力证明，也是他对中国游记传统一个实实在在的拓展：在他之前，很少有人如此自觉地、系统地，把书斋与山川、版本与旅程，写成同一件事。


\medskip\hrule\medskip

\subsection{六、时间的纵深：一座城，他写了四十年}


黄裳留给后来编者最珍贵的结构性遗产，是\textbf{"同地重访"}——同一个地方，他在相隔数十年的不同时代里写过不止一次。


南京是最典型的例子。从 1942 年的《白门秋柳》,到 1946—47 年战后重访的密集书写，到 1949 年时代转折的现场记录，再到 1979 年文革之后的重游追忆——同一座城，他断断续续写了将近四十年。鸡鸣寺更是把这种纵深浓缩进了篇章内部：1946 年的《鸡鸣寺》写豁蒙楼被征用、写后湖紫金山、写胭脂井的旧典；而约 1980 年的《重过鸡鸣寺》，则把陈后主的亡国、十七世纪余怀父子的《金陵览古》、与重游当下,叠成三重时空的回声。把这样两篇并置阅读，你看到的就不只是一座寺的今昔，而是\textbf{一个人的眼光如何随四十年的时代变迁而沉淀、加厚、转调}。1946 年那个愤激于"劫收"乱象的青年记者，与 1980 年那个历尽沧桑、笔调转为低沉的老人，在同一道鸡鸣寺的台阶上完成了一场跨越四十年的对话。


这正是任何一部按原书名、原出版年代排列的旧选本所无法呈现的东西。单行本各自为政时，这些"同地重访"之作散落在不同集子里，彼此不相闻问；唯有把它们抽出来、按地点与年代重新并置，那条贯穿其间的、属于一个人也属于一个时代的时间轴线，才会显形。\textbf{黄裳的全部行旅文字，合起来其实是一部以地点为经、以四十年为纬的私人当代史。} 重编这部选本的意义，一大半就在于把这条被原有书目结构遮蔽了的纵深，重新打捞出来——这既是编辑的工作，也是批评的发现。


\medskip\hrule\medskip

\subsection{结语：一种行将失传的眼光}


把以上各节收拢，黄裳文史游记的特色可以归结为一句话：\textbf{他以一副贯通古今、以古照今的眼光，把行旅、读史与访书熔于一炉，在风景里读人，在风土里读史，在古今的接缝处安身。} 历史品味给了他深度，文章趣味给了他分寸，风土人情给了他体温，而贯穿这一切的，是那种始终把眼前之地当作历史现场来"核对"的自觉。


放在现代中国散文史的坐标里，黄裳是一个承前启后、却又难有来者的存在。他上承晚明小品与周氏兄弟——取周作人的博雅闲适，又取鲁迅的清醒骨力；他身历二十世纪中叶的易代之变，把一代知识人在历史关头的体验，悄悄织进了对三百年前那场易代的回望之中。他既是最后一批能把古书、古迹、古人当作活物来对话的旧式读书人，又是少数能以现代记者之眼为历史作现场见证的新式知识分子。这两种身份在他身上的叠合,几乎是不可复制的——它需要一个人同时拥有藏书家的学问、史家的清醒、文章家的趣味，以及一个亲历大时代转折者的切肤之感。


今天我们重读黄裳，不仅是重读一批好文章，更是在重新学习一种\textbf{看待土地与时间的方式}：每一处风景背后都立着人，每一个人背后都拖着史，而每一段历史，又都在叩问今天的我们将如何选择。这种眼光正在变得稀有。把黄裳散落各处的行旅文字重新编为一集、并为它写下这篇品论，便是想让这副贯通古今的眼光，在它所注视过的那些城与山、那些书与人之间，继续看下去。


\medskip\hrule\medskip

\subsubsection{附录一　本文征引篇目举要}


\begin{quote}\small\itshape
下列为正文论述所倚重的篇目及其要点，依本项目数据库（\texttt{data/huangshang\_essays.json}）整理，供复核与延伸阅读。年份、出处凡标"约""待核"者，须以原书或《黄裳集》最终核定；详细标注（确定／待核／推测）见数据库。所列要点为原创转述，非原文移录。
\end{quote}

\textbf{历史品味·明清易代诸篇}

\begin{itemize}
\item 《石巢园》（《金陵五记》，约 1979）——阮大铖宅园，计成设计，诨名"裤子裆"；以园之美与主之降清对照，写气节之失。
\item 《扫叶楼》（《金陵五记》，约 1979）——清凉山龚贤（1618—1689）故居画室；龚贤着僧服作"扫叶"状以示不事清廷。
\item 《关于吴梅村》（《银鱼集》）——吴伟业贰臣自悔，及《圆圆曲》背后的道德煎熬。
\item 《咏怀堂》（《银鱼集》）——冒辟疆斋号，牵出《影梅庵忆语》与董小宛；注意与《金陵杂记》子目"咏怀堂诗"相区分。
\item 《余淡心与金陵》（《银鱼集》）——余怀游宦秦淮与《板桥杂记》的底本意义。
\item 《绝代的散文家——张宗子》（《银鱼集》）——推许张岱山水小品与故国之思。
\item 《王介甫与金陵》《南唐二陵》（《金陵五记》，约 1979）——上溯北宋与南唐，证其史识之广。
\end{itemize}


\textbf{怀古与见证·1946—1949 金陵现场}

\begin{itemize}
\item 《老虎桥边看"知堂"》（《金陵五记》，1946）——记者身份探看狱中周作人，审奸语境下的人物现场。
\item 《泽存书库》（《金陵五记》，1946）——陈群所建藏书库（名取《礼记》"手泽存焉"），书史与汪伪兴亡交织。
\item 《司徒雷登夹着皮包走了以后》（《金陵五记》，1949）——美国势力退出与国府垮台，与毛泽东同年文章互文。
\item 《解放后看江南》（《金陵五记》，1949）——解放之欢欣与劫后之疮痍并写。
\end{itemize}


\textbf{文章趣味·闲笔、抄书与反讽}

\begin{itemize}
\item 《昆明杂记》（《锦帆集外》，约 1948）——购南明永历野史，以鲁迅"抄古书"笔法结撰。
\item 《采石·当涂·青山》——引陆游《入蜀记》、李白《望天门山》、范传正迁葬青山，三代文献串于一江。
\item 《S.O.S.》《关于"翻译官"》（《关于美国兵》）——后勤荒诞与译员身份的反讽与自嘲（"姨太太""二毛子"）。
\item 《重过鸡鸣寺》（约 1980）——陈后主亡国、余怀父子《金陵览古》、1980 重游三重时空对话；豁蒙楼系张之洞为纪念门生杨锐而建。
\end{itemize}


\textbf{风土人情·以小见史}

\begin{itemize}
\item 《茶馆》（约 1944—1946）——广元／成都／重庆茶馆生态差异，警察与挑夫同座的社会横切面。
\item 《为美国兵活着的人们》（《关于美国兵》）——依附美军经济的底层群像与阶级观察。
\item 《"美人肝"》（《金陵五记》，1946）——南京名馔（鸭胰），马祥兴清真馆，"京苏菜四大名菜"之一。
\end{itemize}


\textbf{三事一体·书旅即行旅}

\begin{itemize}
\item 《西泠访书记》（《榆下说书》）——杭州孤山书肆访书，访书即游山。
\item 《琉璃厂》（见《黄裳自选集》）——来薰阁（始创清咸丰）、1925 年陈济川接掌，访书即城市考古。
\end{itemize}


\textbf{时间纵深·同地重访}

\begin{itemize}
\item 《鸡鸣寺》（1946）／《重过鸡鸣寺》（约 1980）；《金陵杂记》（1947，二十二则子目）与 1979 年《白鹭洲公园补记》《梅园》《莫愁湖》诸篇——同地、隔代、并读。
\end{itemize}


\subsubsection{附录二　与本项目其他文献的关系}


\begin{itemize}
\item 本文的论点与《选编原则》（\texttt{notes/editorial\_principles.md}）"黄裳游记＝地点触发的文史随笔"的根本判断一脉相承，是其在批评层面的展开。
\item 本文第六节"同地重访"与八辑结构（\texttt{outline/new\_anthology\_outline.md}）、各辑导读（\texttt{outline/section\_introductions.md}）相互印证，可作为全书总序（\texttt{outline/volume\_preface\_draft.md}）的理论后盾。
\item 所有篇目事实以数据库为单一事实来源；本文不另立新说，凡超出数据库已有标注的论断，均属批评性阐释而非史实新证。
\end{itemize}


\subsubsection{附录三起　完整篇目资料（据数据库汇编）}


\begin{quote}\small\itshape
为使本文成为一份可独立查考、资料齐备的研究文献，以下数节把本项目迄今\textbf{搜集整理的全部篇目资料}一并收入论文成稿（PDF）： - \textbf{附录三　黄裳行旅篇目总表（全五十九篇，含原创摘要）}——逐篇列出篇名、原收书、年代、地点／地域、取舍（必收／可选／资料库）与置信（确定／待核／推测）标注，及每篇 ≤200 字的原创内容摘要。 - \textbf{附录四　收录书目（书锚）与版本}——据以采辑的原始结集与全集本（《锦帆集》《锦帆集外》《关于美国兵》《金陵五记》《音尘集》等十部）及其版本要点。 - \textbf{附录五　黄裳行旅年表}——按编年汇集可考的行旅与写作节点，供总序附录与读者参照。 - \textbf{附录六　待核清单与已知缺口}——尚待以原书或图书馆书目坐实的条目，以及目前篇目覆盖的薄弱处，透明列出以便续补。 上述四节由本项目数据库 \texttt{data/huangshang\_essays.json} 与 \texttt{notes/timeline\_draft.md} 经 \texttt{data/build\_critique\_pdf.py} \textbf{自动汇编}而成，与全书《编选方案》PDF 同源同步——既保证已整理的宝贵资料完整保留、不致散失，又避免与单一事实来源各自维护而产生歧异。各条置信与来源标注一仍其旧，读者可径据其复核或增补。
\end{quote}

\clearpage

\section{附录三　黄裳行旅篇目总表（全五十九篇，含原创摘要）}
\noindent\small 全表 59 篇，按拟定辑次排列。每篇标注〔编号｜取舍｜置信〕。
摘要均为原创转述（≤200 字），\textbf{非原文复制}；篇名、年代、出处凡未坐实者标「待核 / 推测」。\normalsize
\medskip
\subsection{第1辑　蜀道与战时中国（9 篇）}
\noindent\textbf{《宝鸡——广元》}\,{\small ［E-014｜★必收｜待核］}\\
{\small\itshape 原收入：锦帆集　｜　年代：约1944—1946　｜　地点：宝鸡（陕）—广元（川），蜀道（蜀中）}\\
日记体纪行，记由陕西宝鸡入川至广元一线蜀道行程：穿越秦岭、越剑门、入川盆地，描绘沿途关隘峰壑与战时人流；亦收入《往事如烟》等选集。第一辑蜀道行旅最核心的地理线索篇。
\par\vspace{5pt}
\noindent\textbf{《成都散记》}\,{\small ［E-015｜★必收｜待核］}\\
{\small\itshape 原收入：锦帆集　｜　年代：约1944—1946　｜　地点：成都（蜀中）}\\
1943—1946年任美军译员期间对成都后方生活的散记：写宽街窄巷、市井风物与战时文化氛围；与《茶馆》（E-036）形成互补，共同呈现抗战时期蜀中都市的生活横切面。《锦帆集》蜀道组核心篇之一。
\par\vspace{5pt}
\noindent\textbf{《过徐州》}\,{\small ［E-016｜○可选｜待核］}\\
{\small\itshape 原收入：锦帆集　｜　年代：约1944—1946　｜　地点：徐州（江北）}\\
1942年冬（腊月初旬）离沪赴蜀途中过徐州：见集市陈列布匹铜器，与同行者（含黄宗江、宋淇等）购买徐州风景明信片寄回上海；是《锦帆集》蜀道诸篇序章，战时离乡漂泊之感跃然。
\par\vspace{5pt}
\noindent\textbf{《江上杂记》}\,{\small ［E-017｜○可选｜待核］}\\
{\small\itshape 原收入：锦帆集　｜　年代：约1944—1946　｜　地点：长江（舟行，具体江段待核）（江南）}\\
舟行江上的杂记，记水路见闻与过境感受；取书名意象'锦帆应是到天涯'之境。确切江段（川江或长江下游）待核，归辑随之调整；是《锦帆集》少数以江河场景为主的行旅篇。
\par\vspace{5pt}
\noindent\textbf{《昆明杂记》}\,{\small ［E-018｜★必收｜待核］}\\
{\small\itshape 原收入：锦帆集外（文学丛刊第九集，文化生活出版社，1948；含五部分）　｜　年代：约1944—1946　｜　地点：昆明（西南）}\\
任译员期间在昆明的见闻杂记，分五部分；在昆明购得载永历帝滇缅流亡的晚明野史，三百年后重走同一片土地，遂以鲁迅'抄古书'笔法写历史笔记式杂记，沧桑感浓；文气近汪曾祺。收《锦帆集外》（文学丛刊第九集）。
\par\vspace{5pt}
\noindent\textbf{《森林·雨季·山头人》}\,{\small ［E-019｜○可选｜待核］}\\
{\small\itshape 原收入：锦帆集（待核）　｜　年代：约1944—1946　｜　地点：西南/滇缅印一带（待核）（西南／海外）}\\
战时行役中关于森林、雨季与边地人群的记述，地理指向待核（可能涉滇缅印战区）。具民族志色彩。
\par\vspace{5pt}
\noindent\textbf{《茶馆》}\,{\small ［E-036｜○可选｜待核］}\\
{\small\itshape 原收入：锦帆集（待核）　｜　年代：约1944—1946　｜　地点：西南（成渝一带，待核）（蜀中／西南）}\\
写西南后方各地茶馆的生态差异与阶层混杂：广元/剑阁茶叶粗砺，成都茶馆松弛，重庆茶馆兼作谈判场——警察与挑夫同座、大学生以茶馆为自修室、生意人在此成交。一室之中浓缩战时社会横截面，风俗志价值极高。
\par\vspace{5pt}
\noindent\textbf{《断片》}\,{\small ［E-055｜○可选｜待核］}\\
{\small\itshape 原收入：锦帆集　｜　年代：约1944—1946　｜　地点：（追忆，非行旅地点）（蜀中／西南）}\\
《锦帆集》第二章，题名'断片'；据研究，此篇追忆少年时期的初恋情感，是全书调性最私密的篇目，与其余蜀道行旅诸篇形成对比；作为全集青春记忆的底色，是理解黄裳早年情感世界的罕见文本。
\par\vspace{5pt}
\noindent\textbf{《音尘》}\,{\small ［E-056｜○可选｜推测］}\\
{\small\itshape 原收入：锦帆集　｜　年代：约1944—1946　｜　地点：西南后方（待核）（蜀中／西南）}\\
《锦帆集》第七章，以'音尘'为题（音讯，兼含离别与消逝之叹）；据研究，写战时音讯阻隔与离别之痛，是全书情感基调最深沉的一篇；《音尘集》（1996年辽宁教育版）即以此篇命名，彰显其在黄裳自我认知中的特殊分量。
\par\vspace{5pt}
\subsection{第2辑　美国兵与异国战争经验（16 篇）}
\noindent\textbf{《叙言》}\,{\small ［E-020｜·资料库｜待核］}\\
{\small\itshape 原收入：关于美国兵　｜　年代：约1947　｜　地点：—（—）}\\
《关于美国兵》全书序言，交代写作缘起与立场。导读价值高。
\par\vspace{5pt}
\noindent\textbf{《他们怎样看中国》}\,{\small ［E-021｜★必收｜待核］}\\
{\small\itshape 原收入：关于美国兵　｜　年代：约1944—1947　｜　地点：西南后方（待核）（西南）}\\
记录美军士兵对中国的具体印象：普通士兵称中国人'聪明高尚但缺乏教育'，芝加哥大学毕业军官称中国人'缺乏组织、散漫成群'；以具体人物的实际言论呈现中美文化认知落差，是《关于美国兵》中中美文化观察最集中的篇章。
\par\vspace{5pt}
\noindent\textbf{《“中国通”》}\,{\small ［E-022｜·资料库｜待核］}\\
{\small\itshape 原收入：关于美国兵　｜　年代：约1944—1947　｜　地点：西南后方（待核）（西南）}\\
写美军中自诩'中国通'的外籍军官或专家：以具体人物的语言习惯与认知偏差呈现文化误读与傲慢，讽刺意味浓；与《他们怎样看中国》（E-021）构成中美文化认知专题的两面。
\par\vspace{5pt}
\noindent\textbf{《人种·职业·人性的大集合》}\,{\small ［E-023｜○可选｜待核］}\\
{\small\itshape 原收入：关于美国兵　｜　年代：约1944—1947　｜　地点：西南后方（待核）（西南）}\\
写美军内部各色人种（黑人、白人、拉丁裔等）、职业（军官/士兵/文职）与人性混杂景观；从社会学角度呈现美军作为一个社会横截面的复杂性，群像式观察，史料感强。
\par\vspace{5pt}
\noindent\textbf{《几个人物》}\,{\small ［E-024｜○可选｜待核］}\\
{\small\itshape 原收入：关于美国兵　｜　年代：约1944—1947　｜　地点：西南后方（待核）（西南）}\\
记战时接触的几位美军人物素描：通过具体个人的性格、言行与经历呈现美军群像；人物速写式，文学性强，与第八辑'人物途中'的写法呼应。具体人物待核原文。
\par\vspace{5pt}
\noindent\textbf{《中国将军怎样应付美国兵》}\,{\small ［E-025｜★必收｜待核］}\\
{\small\itshape 原收入：关于美国兵　｜　年代：约1944—1947　｜　地点：西南后方（待核）（西南）}\\
写中国将军应对美军联合指挥的方式：以官场智慧周旋于美方僵化指挥与本方利益之间，现场政治与人情交织；以具体军官行为呈现中美联军权力结构中的张力与博弈。
\par\vspace{5pt}
\noindent\textbf{《美国兵与女人》}\,{\small ［E-026｜★必收｜待核］}\\
{\small\itshape 原收入：关于美国兵　｜　年代：约1944—1947　｜　地点：西南后方（待核）（西南）}\\
写战时美军与中国社会接触中关于女性的一面：美军基地周边的娱乐场所、中美战时婚恋关系与性别风气变迁；纪实性强，是《关于美国兵》最具社会学价值的一章，亦见《黄裳自选集》。
\par\vspace{5pt}
\noindent\textbf{《“伟大”的S.O.S.》}\,{\small ［E-027｜·资料库｜待核］}\\
{\small\itshape 原收入：关于美国兵　｜　年代：约1944—1947　｜　地点：西南后方（待核）（西南）}\\
S.O.S.即美军后勤勤务部（Services of Supply），以反讽笔调写其'伟大'运作：繁复补给程序、物资浪费与官僚机器荒诞；物质批判与战争荒诞感合一，标题引号即反讽所在。
\par\vspace{5pt}
\noindent\textbf{《前线景象》}\,{\small ［E-028｜○可选｜待核］}\\
{\small\itshape 原收入：关于美国兵　｜　年代：约1944—1947　｜　地点：前线（待核，疑滇西/缅北）（西南／海外）}\\
记前线见闻：战壕、废墟、炮火下的日常与死亡气息；具体战区待核（疑滇西或缅北反攻一线），是《关于美国兵》中从后方视角转向真实战场的一章，现场感最强。
\par\vspace{5pt}
\noindent\textbf{《咖啡与战斗力》}\,{\small ［E-029｜·资料库｜待核］}\\
{\small\itshape 原收入：关于美国兵　｜　年代：约1944—1947　｜　地点：西南后方（待核）（西南）}\\
以咖啡这一美军标配物资为切口写后勤与战斗力：充足补给、规律饮食与高昂士气，与中国军队物质匮乏形成对照；借小物件见战争逻辑与国力差异，角度新颖独到。
\par\vspace{5pt}
\noindent\textbf{《种种惊异》}\,{\small ［E-030｜○可选｜待核］}\\
{\small\itshape 原收入：关于美国兵　｜　年代：约1944—1947　｜　地点：西南后方（待核）（西南）}\\
记中美接触中的种种文化惊异：饮食习惯、礼仪观念、价值判断的错位与误解；以具体场景呈现文化落差，宜与《他们怎样看中国》（E-021）对读，共构第二辑文化碰撞主题。
\par\vspace{5pt}
\noindent\textbf{《军中文化》}\,{\small ［E-031｜·资料库｜待核］}\\
{\small\itshape 原收入：关于美国兵　｜　年代：约1944—1947　｜　地点：西南后方（待核）（西南）}\\
写美军中丰富的文化生活与阅读资源：美军随军携带《生活》《时代》《展望》《纽约客》等杂志，驻地自办军中报纸，书架备有奥尼尔戏剧集、狄更斯小说等文学经典；与中国战时文化资源的匮乏形成强烈对照。
\par\vspace{5pt}
\noindent\textbf{《汽车团》}\,{\small ［E-032｜·资料库｜待核］}\\
{\small\itshape 原收入：关于美国兵　｜　年代：约1944—1947　｜　地点：西南后方（待核，疑滇缅公路沿线）（西南）}\\
写美军汽车团沿滇缅公路长途卡车运输的艰险：历经'二十四道拐'险途与西南山区路况，驾驶员以机器维系后方命脉；以具体运输生态呈现抗战后勤战争的真实面貌，是《关于美国兵》中场景感最强的一章。
\par\vspace{5pt}
\noindent\textbf{《为美国兵活着的人们》}\,{\small ［E-033｜★必收｜待核］}\\
{\small\itshape 原收入：关于美国兵　｜　年代：约1944—1947　｜　地点：西南后方（待核）（西南）}\\
写战时围绕美军基地谋生的中国底层人群：街头摊贩、洗衣女工、酒吧女招待、翻译助手……生计与尊严皆系于美军的存在。以底层视角呈现战时经济依附性，社会批判意味浓烈，是《关于美国兵》中阶级观察最深刻的一章。
\par\vspace{5pt}
\noindent\textbf{《关于“翻译官”》}\,{\small ［E-034｜★必收｜待核］}\\
{\small\itshape 原收入：关于美国兵　｜　年代：约1944—1947　｜　地点：西南后方（待核）（西南）}\\
以作者本人译员经历为视角，辨析'翻译官'称呼背后的身份尴尬：外事局称'X级译员'，远征军称'Interpreting Officer'，民间戏称'姨太太''二毛子'，折射中美权力关系在日常语言中的投影；夫子自道，是第二辑最贴近作者战时身份的收束之篇。
\par\vspace{5pt}
\noindent\textbf{《印度小夜曲》}\,{\small ［E-060｜★必收｜待核］}\\
{\small\itshape 原收入：音尘集（辽宁教育出版社，书趣文丛第三辑，1996）　｜　年代：约1944—1946　｜　地点：印度（加尔各答、菩提加雅等）（海外）}\\
黄裳1944—1946年随中美混合部队赴印度期间的行旅长篇，分八章：序曲、雨季、加尔各答、菩提加雅、群莺乱飞、弄蛇者、女人们、森林雨季。以南亚异域风土与战时情境为主要内容，是《关于美国兵》之外战时海外经验的另一维度。
\par\vspace{5pt}
\subsection{第3辑　金陵旧梦（24 篇）}
\noindent\textbf{《白门秋柳》}\,{\small ［E-001｜★必收｜确定］}\\
{\small\itshape 原收入：锦帆集（1946，写于1942；后收入《金陵五记》）　｜　年代：1942　｜　地点：南京（白门）（金陵）}\\
1942年冬黄裳与黄宗江、宋淇（宋希）、周杲良四人离沪赴川，途经汪伪'首都'南京；描绘沦陷南京的衰飒——鸡鸣寺、随园、扫叶楼凋落，书店中《人间味》杂志与'今日的江南的无声的悲哀'；称之为'揉碎了的美'的战时见证。
\par\vspace{5pt}
\noindent\textbf{《旅京随笔》}\,{\small ［E-002｜○可选｜待核］}\\
{\small\itshape 原收入：金陵五记　｜　年代：1946　｜　地点：南京（金陵）}\\
1946年黄裳以记者身份赴南京，目睹国民政府'劫收'乱象冲击民生文化；乘机游览名胜、探访旧日友人与藏书家。'旅京随笔'是《金陵五记》1946年编总称，下含鸡鸣寺、泽存书库、盔山精舍、老虎桥知堂、美人肝共五篇。
\par\vspace{5pt}
\noindent\textbf{《鸡鸣寺》}\,{\small ［E-003｜★必收｜待核］}\\
{\small\itshape 原收入：金陵五记　｜　年代：1946　｜　地点：南京·鸡鸣寺（金陵）}\\
1946年黄裳以记者身份访鸡鸣寺，欲于豁蒙楼写稿，发现楼已被'防空司令部'征用，只得在大殿内喝茶，眺望后湖紫金山；记陈后主'胭脂井'典故（三人同束以绳引出）与南朝兴亡之感。与1979年《重过鸡鸣寺》构成同地重访。
\par\vspace{5pt}
\noindent\textbf{《关于“泽存书库”》}\,{\small ［E-004｜○可选｜待核］}\\
{\small\itshape 原收入：金陵五记　｜　年代：1946　｜　地点：南京（金陵）}\\
记汪伪官员陈群在南京所建'泽存书库'（取《礼记》'手泽存焉'之义）：藏书40万册、善本4400余部，含汪精卫旧藏；战后被接收，善本随蒋携台。访书过程中书史与汪伪政权兴亡交织，史料价值极高。
\par\vspace{5pt}
\noindent\textbf{《访“盔山精舍”》}\,{\small ［E-005｜○可选｜待核］}\\
{\small\itshape 原收入：金陵五记　｜　年代：1946　｜　地点：南京（金陵）}\\
访南京'盔山精舍'的记游兼记人之作。所涉人物与建置待核。
\par\vspace{5pt}
\noindent\textbf{《老虎桥边看“知堂”》}\,{\small ［E-006｜★必收｜待核］}\\
{\small\itshape 原收入：金陵五记　｜　年代：1946　｜　地点：南京·老虎桥监狱（金陵）}\\
1946年黄裳以《文汇报》记者身份赴南京，专赴老虎桥模范监狱探看周作人（知堂）；经延龄巷、国府路找到监狱，申请获准后会面，完成于最终宣判前。长篇报道性质，是战后审奸语境下的人物现场记录，史料价值极高。
\par\vspace{5pt}
\noindent\textbf{《金陵杂记》}\,{\small ［E-007｜★必收｜待核］}\\
{\small\itshape 原收入：金陵五记　｜　年代：1947　｜　地点：南京（多处）（金陵）}\\
1947年系列短札，下含小序及半山寺与谢公墩、石观音寺、周处读书台、绛云书卷美人图、柳如是、快园、随园、梅园、后湖、小虹桥、咏怀堂诗、梅花山、莫愁湖、明太祖与徐达、桃花扇底看南朝、马湘兰、记者生涯、豁蒙楼上看浓春、燕子矶、白鹭洲、鸡鹅巷与裤子裆等二十二则。
\par\vspace{5pt}
\noindent\textbf{《解放后看江南》}\,{\small ［E-008｜★必收｜待核］}\\
{\small\itshape 原收入：金陵五记　｜　年代：1949　｜　地点：江南／南京（江南）}\\
1949年秋南京解放，黄裳亲历人民迎接解放的欢欣；同时目睹满目疮痍、民生凋蔽——劫后时代的两面。历史巨变的现场见证，兼具时代兴奋与直面现实的张力，是《金陵五记》1949年编核心篇。
\par\vspace{5pt}
\noindent\textbf{《司徒雷登夹着皮包走了以后》}\,{\small ［E-009｜○可选｜待核］}\\
{\small\itshape 原收入：金陵五记　｜　年代：1949　｜　地点：南京（金陵）}\\
1949年以美国驻华大使司徒雷登携皮包黯然离华为楔子，写南京政局剧变中美国势力退出与国民政府垮台的亲历观察；与毛泽东同年《别了，司徒雷登》形成历史互文。时评性强，史料价值可观。
\par\vspace{5pt}
\noindent\textbf{《白下书简》}\,{\small ［E-010｜○可选｜待核］}\\
{\small\itshape 原收入：金陵五记　｜　年代：1979　｜　地点：南京（白下）（金陵）}\\
1979年以书简体重写南京，收十一篇；格调低沉，距早年金陵记忆逾三十年，文革后重访兼含追忆旧迹与今昔对照。'白下'为南京历史别称。是《金陵五记》1979年编五类文字之一。
\par\vspace{5pt}
\noindent\textbf{《秦淮拾梦记》}\,{\small ［E-011｜★必收｜待核］}\\
{\small\itshape 原收入：金陵五记　｜　年代：1979　｜　地点：南京·秦淮（金陵）}\\
1979年重游秦淮，以余怀《板桥杂记》为引，追溯秦淮旧影；涉及柳如是、陈圆圆、余淡心等明末清初人物；是1979年南京访古组最具文学性的一篇，网络有节选流传（须以原书核全文）。
\par\vspace{5pt}
\noindent\textbf{《重过鸡鸣寺》}\,{\small ［E-012｜★必收｜待核］}\\
{\small\itshape 原收入：金陵五记　｜　年代：1979　｜　地点：南京·鸡鸣寺（金陵）}\\
约1980年写的重访记，构建三重时空对话：陈后主亡国（胭脂井典故）、十七世纪余怀父子文字（《金陵览古》）、1980年重游；豁蒙楼本为张之洞纪念门生杨锐所建，几度易主；是全书'同地重访'结构最完整的一篇。
\par\vspace{5pt}
\noindent\textbf{《南唐二陵》}\,{\small ［E-013｜○可选｜待核］}\\
{\small\itshape 原收入：金陵五记　｜　年代：1979　｜　地点：南京·南唐二陵（金陵）}\\
1979年访南京祖堂山南唐二陵（先主李{\fb 昪}钦陵、中主李{\fb 璟}顺陵），以词帝李煜亡国为历史背景；追溯南唐短暂王朝（仅39年）的文艺辉煌与政治覆灭，是金陵兴亡感最浓烈的一篇怀古文字。
\par\vspace{5pt}
\noindent\textbf{《「美人肝」》}\,{\small ［E-044｜○可选｜待核］}\\
{\small\itshape 原收入：金陵五记　｜　年代：1946　｜　地点：南京（金陵）}\\
写南京名馔'美人肝'（鸭胰脏）：此菜与松鼠鱼、蛋烧卖、凤尾虾并称'京苏菜四大名菜'，品味于清真老字号马祥兴馆；借饮食掌故钩沉金陵饮食文化传统，文气轻灵，兼具考证趣味与口腹之乐，是《金陵五记》1946年组中风格最轻松的一篇。
\par\vspace{5pt}
\noindent\textbf{《社会科学者眼里的南京》}\,{\small ［E-045｜○可选｜待核］}\\
{\small\itshape 原收入：金陵五记　｜　年代：1949　｜　地点：南京（金陵）}\\
1949年以社会观察者视角审视南京政治、经济与社会变动。内容待核。与同年《解放后看江南》《司徒雷登夹着皮包走了以后》同属1949年南京见证书写。
\par\vspace{5pt}
\noindent\textbf{《不再做花瓶——记中国农业机械公司南京工厂》}\,{\small ［E-046｜○可选｜待核］}\\
{\small\itshape 原收入：金陵五记　｜　年代：1949　｜　地点：南京（金陵）}\\
1949年记南京某工厂（中国农业机械公司）的纪实文字，属解放初期社会面貌观察。内容与文学质量待核。
\par\vspace{5pt}
\noindent\textbf{《白鹭洲公园补记》}\,{\small ［E-047｜○可选｜待核］}\\
{\small\itshape 原收入：金陵五记　｜　年代：1979　｜　地点：南京·白鹭洲（金陵）}\\
1979年对白鹭洲公园的补记。白鹭洲为南京历史园林，有秦淮文化脉络；此篇与《金陵杂记》1947年白鹭洲条构成同地重访对照。内容待核。
\par\vspace{5pt}
\noindent\textbf{《天王府的西花园》}\,{\small ［E-048｜★必收｜待核］}\\
{\small\itshape 原收入：金陵五记　｜　年代：1979　｜　地点：南京·天王府（太平天国旧址）（金陵）}\\
1979年访太平天国天王府遗址之西花园（今南京总统府区域）：以历史层积视角呈现一地三叠——太平天国洪秀全宫室、民国国民政府、当代公园；黄裳追溯晚清历史地图与今昔之变，金陵近代史迹密度最高、时间跨度最宽的一篇。
\par\vspace{5pt}
\noindent\textbf{《梅园（1979版）》}\,{\small ［E-049｜○可选｜待核］}\\
{\small\itshape 原收入：金陵五记　｜　年代：1979　｜　地点：南京·梅园（金陵）}\\
1979年专题写梅园新村（1946年国共谈判中共代表团驻地，周恩来曾居此）：地点今已辟为纪念馆，以亲历者视角追忆三十年前谈判史迹与政治氛围；与《金陵杂记》1947年梅园子目构成时间纵深，重访结构清晰。
\par\vspace{5pt}
\noindent\textbf{《王介甫与金陵》}\,{\small ［E-050｜★必收｜待核］}\\
{\small\itshape 原收入：金陵五记　｜　年代：1979　｜　地点：南京（半山园等王安石相关遗迹）（金陵）}\\
写王安石（字介甫）晚年退居金陵的行迹与精神面貌：'骑驴粗服往来钟山'；追溯其《桂枝香·金陵怀古》被后人称为绝唱（'诸公寄调桂枝香者三十余家，惟王介甫为绝唱'）；政治家与词人两面在金陵山川中交叠，史论与地理感并茂。
\par\vspace{5pt}
\noindent\textbf{《莫愁湖（1979版）》}\,{\small ［E-051｜○可选｜待核］}\\
{\small\itshape 原收入：金陵五记　｜　年代：1979　｜　地点：南京·莫愁湖（金陵）}\\
1979年专题写莫愁湖：钩沉'莫愁女'传说（南朝民歌《莫愁乐》主人公）与'胜棋楼'朱元璋赐名典故；与《金陵杂记》1947年莫愁湖条构成三十年重访对照，展现金陵历史地标在现代的延续与变迁。须核定是否独立成篇。
\par\vspace{5pt}
\noindent\textbf{《石巢园》}\,{\small ［E-052｜○可选｜待核］}\\
{\small\itshape 原收入：金陵五记　｜　年代：1979　｜　地点：南京·石巢园（金陵）}\\
1979年访南京石巢园——明末清初人物阮大铖的宅园（由造园大家计成参与设计）；阮大铖以'库司坊'（民间蔑称'裤子裆'）筑此园，后降清变节；黄裳以园林精美与园主人格的反差呈现历史道德的复杂性，文史批判力度强。
\par\vspace{5pt}
\noindent\textbf{《扫叶楼》}\,{\small ［E-053｜○可选｜待核］}\\
{\small\itshape 原收入：金陵五记　｜　年代：1979　｜　地点：南京·扫叶楼（清凉山）（金陵）}\\
1979年访南京清凉山扫叶楼：楼以清初画家龚贤（1618—1689）旧居兼画室得名，龚贤曾穿僧服作扫叶状以示不事清廷；黄裳追溯金石画史与气节人格，书画史与金陵山川地脉交织，是第三辑中画史掌故最丰富的一篇。
\par\vspace{5pt}
\noindent\textbf{《献花岩》}\,{\small ［E-054｜○可选｜待核］}\\
{\small\itshape 原收入：金陵五记　｜　年代：1979　｜　地点：南京郊外·献花岩（具体地点待核）（金陵）}\\
1979年访献花岩（南京近郊某山川古迹）。具体位置与内容待核，可能涉山川探胜与历史追想。
\par\vspace{5pt}
\subsection{第4辑　江南行旅（3 篇）}
\noindent\textbf{《采石·当涂·青山》}\,{\small ［E-035｜○可选｜待核］}\\
{\small\itshape 原收入：锦帆集（待核）　｜　年代：待核　｜　地点：采石矶、当涂、青山（皖南／长江沿线）（江南）}\\
沿长江访采石矶、当涂、青山李白墓一线：引陆游《入蜀记》勾勒采石矶山水；述李白初赴当涂所赋《望天门山》；追溯范传正将李白迁葬青山始末。山水、史迹、人物三合一，是第四辑江南行旅的典范之作。
\par\vspace{5pt}
\noindent\textbf{《海滨消夏记》}\,{\small ［E-039｜○可选｜待核］}\\
{\small\itshape 原收入：银鱼集（后并入《黄裳集·创作卷V》）　｜　年代：待核　｜　地点：海滨（具体地待核，疑青岛/北戴河一带）（待核）}\\
海滨（疑青岛或北戴河）避暑期间的行旅散记，旁及散文美学漫想（'科学著作也可成美文'），文风轻松闲适；属黄裳较少见的休闲题材，可丰富全书行旅调性的层次，与战时纪行形成对照。
\par\vspace{5pt}
\noindent\textbf{《杭州杂记》}\,{\small ［E-043｜○可选｜待核］}\\
{\small\itshape 原收入：花步集（待核）　｜　年代：待核　｜　地点：杭州（江南）}\\
写杭州（西湖一带）的城市游记，已见于《黄裳散文选集》第一辑记游篇目，篇目真实性较可靠；具体内容与原书归属（疑《花步集》）仍待核，须以原书确认篇名、文字与归集。
\par\vspace{5pt}
\subsection{第5辑　北平与旧都余韵（1 篇）}
\noindent\textbf{《琉璃厂》}\,{\small ［E-038｜★必收｜待核］}\\
{\small\itshape 原收入：待核（见《黄裳自选集》第四辑；疑出书话诸集）　｜　年代：待核　｜　地点：北京·琉璃厂（北平）}\\
写北京琉璃厂旧书古玩街的访书见闻：从清咸丰年间始创的来薰阁，到1925年陈济川接掌后旧书业的领军地位；书肆掌故与北平城市记忆交织，在书话与行旅之间打通'书旅'主题，是第五辑故都题材与第七辑书肆题材的双重核心篇。
\par\vspace{5pt}
\subsection{第7辑　书肆、旧书与纸上行旅（1 篇）}
\noindent\textbf{《西泠访书记》}\,{\small ［E-037｜★必收｜待核］}\\
{\small\itshape 原收入：榆下说书（后并入《黄裳集·创作卷V》）　｜　年代：待核　｜　地点：杭州·西泠（西湖孤山一带）（江南）}\\
记在杭州西湖孤山一带（西泠）寻访旧书铺的经历：书肆散布于园林湖山之间，访书即是行旅，翻检中偶遇珍本；书事、人事与西湖景致交织，是'旅行即读书'主题的最佳演示，亦是第七辑书肆旅行专辑的核心篇。
\par\vspace{5pt}
\subsection{第8辑　人物途中（5 篇）}
\noindent\textbf{《余淡心与金陵》}\,{\small ［E-040｜★必收｜待核］}\\
{\small\itshape 原收入：银鱼集（后并入《黄裳集·创作卷V》）　｜　年代：待核　｜　地点：南京（秦淮）（金陵）}\\
以明末清初文人余怀（字淡心，著《板桥杂记》）游宦金陵的行迹为线索：钩沉其在秦淮河畔的交游与写作背景，追究《板桥杂记》的底本意义；以人物串联金陵文学史，是'以人写地'写法的典范，与E-057互补，建议第三辑/第八辑各取其一。
\par\vspace{5pt}
\noindent\textbf{《绝代的散文家——张宗子》}\,{\small ［E-041｜○可选｜待核］}\\
{\small\itshape 原收入：银鱼集（后并入《黄裳集·创作卷V》）　｜　年代：待核　｜　地点：—（涉绍兴/杭州等张岱行迹）（江南）}\\
以'绝代的散文家'高度评价明末张岱（字宗子）：细析其山水小品的感官精准与故国之思，以《陶庵梦忆》《西湖梦寻》为例；黄裳自身趣味与张岱遥相呼应，是理解其文学观的重要参照。属论人之作，行旅成分弱，列可选。
\par\vspace{5pt}
\noindent\textbf{《关于余淡心》}\,{\small ［E-057｜○可选｜待核］}\\
{\small\itshape 原收入：银鱼集（后并入《黄裳集·创作卷V》）　｜　年代：待核　｜　地点：—（涉南京秦淮等余怀相关地点）（金陵）}\\
关于明末清初文人余怀（字淡心，著《板桥杂记》）的评论性记人文字，与E-040《余淡心与金陵》互见。前者以金陵地点为主，此篇侧重人物评述。内容待核。
\par\vspace{5pt}
\noindent\textbf{《关于吴梅村》}\,{\small ［E-058｜○可选｜待核］}\\
{\small\itshape 原收入：银鱼集（后并入《黄裳集·创作卷V》）　｜　年代：待核　｜　地点：—（涉江南，吴伟业相关行迹）（江南）}\\
写清初诗人吴伟业（字骏公，号梅村）的贰臣处境：以诗文成就与仕清后的自悔心理为核心，追溯其《圆圆曲》背后的政治遭遇与道德煎熬；黄裳对晚明清初气节问题的持续关注在此集中体现，是第八辑人物书写中思想张力最强的一篇。
\par\vspace{5pt}
\noindent\textbf{《咏怀堂》}\,{\small ［E-059｜○可选｜待核］}\\
{\small\itshape 原收入：银鱼集（后并入《黄裳集·创作卷V》）　｜　年代：待核　｜　地点：—（涉江苏如皋或南京等冒辟疆相关地点，待核）（江南）}\\
写明末四公子之一冒辟疆（号巢民）的斋号'咏怀堂'：追溯冒辟疆与秦淮名妓董小宛的传奇情事（见《影梅庵忆语》）及其不仕清廷的气节；园林、书斋与人格互照，是江南士大夫文化书写的典型个案。注意与E-007《金陵杂记》子目'咏怀堂诗'（1947年南京所见诗集）区分。
\par\vspace{5pt}

\section{附录四　收录书目（书锚）与版本}
\noindent\small 本项目据以采辑的原始结集与全集本；版权均未失效，仅作版本核校之用。\normalsize\medskip
\noindent\textbf{《锦帆集》}\,{\small（中华书局，1946，中华文艺丛刊）}\\
作者第一部散文集；书名取李商隐（义山）'锦帆应是到天涯'句。丛刊归属已多源印证（确定）。章序：ch1序、ch2断片、ch3白门秋柳、ch4过徐州、ch5宝鸡——广元、ch6成都散记、ch7音尘、ch8江上杂记、ch9'江湖'后记、ch10后记。多记抗战后期西南/蜀道行旅。注意：昆明杂记不在此集，在《锦帆集外》。
\\{\small\itshape 来源：多源印证（豆瓣'锦帆集 锦帆集外 关于美国兵'条目、人民日报副刊、百度百科）；丛刊名'中华文艺丛刊'多源确认。}
\par\vspace{5pt}
\noindent\textbf{《锦帆集外》}\,{\small（文化生活出版社，1948，文学丛刊第九集（巴金主编，多源印证））}\\
《锦帆集》集外文字结集，列巴金主持的'文学丛刊第九集'（多源印证）。核心篇目：《昆明杂记》（五部分；E-018）。完整篇目尚待录入。
\\{\small\itshape 来源：网络检索（豆瓣/微信读书；多处提及巴金文学丛刊第九集）；昆明杂记在此集已多源印证。}
\par\vspace{5pt}
\noindent\textbf{《音尘集》}\,{\small（辽宁教育出版社，1996，书趣文丛第三辑）}\\
含三部分合订：（1）《锦帆集》重印；（2）《印度小夜曲》（新增，8章：序曲/雨季/加尔各答/菩提加雅/群莺乱飞/弄蛇者/女人们/森林雨季；写1944—1946年印度行役）；（3）《关于美国兵》重印。《印度小夜曲》是目前能确认的印度行旅专题长篇，为本项目重要新增来源。
\\{\small\itshape 来源：网络检索（多源提及'音尘集'辽宁教育版，书趣文丛第三辑，1996）；年份与出版社待实物核。}
\par\vspace{5pt}
\noindent\textbf{《关于美国兵》}\,{\small（上海出版公司，1947，周报丛书第一辑）}\\
作者抗战末随美军任译员期间的现场观察。已得完整目录（15篇，见 essays 中 E-020..E-034）。
\\{\small\itshape 来源：多源印证（豆瓣条目、综述文章）。}
\par\vspace{5pt}
\noindent\textbf{《金陵五记》}\,{\small（金陵书画社，1982）}\\
汇集1942—1979年写南京文字，分五个时间层。1982年6月初版，260页，统一书号10234-004，定价¥0.85（金陵书画社版，已多源确认）。另有后出江苏古籍出版社版（ISBN 9787806431757）。旧版录'江苏人民出版社'为误，已更正。全书核心来源。
\\{\small\itshape 来源：网络检索（百度百科'金陵五记'、微信读书/中图网目录、商务'流金文丛'重印页）；金陵书画社1982版书号统一书号10234-004，多源确认。}
\par\vspace{5pt}
\noindent\textbf{《花步集》}\,{\small（生活·读书·新知三联书店，1982）}\\
后并入《黄裳集·创作卷IV》（与《金陵五记》《晚春的行旅》同册）。篇目待录入。
\\{\small\itshape 来源：网络检索（图书馆书目、豆瓣《黄裳集》系列）；出版社/年份单源，待核。}
\par\vspace{5pt}
\noindent\textbf{《晚春的行旅》}\,{\small（待核（疑山西人民出版社，未证实），待核（疑1980年代初））}\\
书名直指'行旅'主题，应为本选本重点采辑对象。后并入《黄裳集·创作卷IV》。篇目与出版信息均待核。
\\{\small\itshape 来源：网络检索（豆瓣《黄裳集》系列）；出版社/年份未坐实。}
\par\vspace{5pt}
\noindent\textbf{《榆下说书 / 银鱼集》}\,{\small（（原版年份待核）后并入《黄裳集·创作卷V》（山东人民出版社），待核）}\\
书话集。《榆下说书》得目录线索：书的故事、古书的作伪、西泠访书记、谈'善本'、谈'题跋'、谈'集部'、谈禁书、书痴等；《银鱼集》：绝代的散文家——张宗子、关于张宗子、关于余淡心、余淡心与金陵、关于吴梅村、咏怀堂、海滨消夏记等。为第七、八辑主要来源。
\\{\small\itshape 来源：网络检索（中图网/苏宁《黄裳集·创作卷V》页、Google Books）。}
\par\vspace{5pt}
\noindent\textbf{《翠墨集 / 榆下杂说》}\,{\small（后并入《黄裳集·创作卷VI》（山东人民出版社），待核）}\\
书话/杂说集，供第六、七辑补采。篇目待录入。
\\{\small\itshape 来源：网络检索（苏宁《黄裳集·创作卷VI》页）。}
\par\vspace{5pt}
\noindent\textbf{《黄裳集（山东人民出版社）》}\,{\small（山东人民出版社，约2017—2022（分批））}\\
目前最全的作品集，分创作卷（约16卷）+古籍研究卷+译文卷+书信卷。建议作为主要底本与核对基准。创作卷I=锦帆集/锦帆集外/关于美国兵；IV=花步集/金陵五记/晚春的行旅；V=榆下说书/银鱼集；VI=翠墨集/榆下杂说。仍在版权保护期内。
\\{\small\itshape 来源：网络检索（豆瓣'黄裳集'丛书页、中图网团购页、苏宁分卷页）。卷次/年份待核。}
\par\vspace{5pt}

\section{附录五　黄裳行旅年表}

\begin{quote}\small\itshape
由 \texttt{data/build\_timeline.py} 自动生成，勿手工修改。 仅含 \texttt{approximate\_year} 有值的篇目；年份均待核，「确定」者另见 \texttt{data/huangshang\_essays.json}。 图例：★ 必收　○ 可选　· 资料库
\end{quote}

\subsection{1942}


\begin{itemize}
\item ★ 【辑3】\textbf{白门秋柳} — 南京（白门）
\end{itemize}


\subsection{约1944—1946}


\begin{itemize}
\item ★ 【辑2】\textbf{印度小夜曲} — 印度（加尔各答、菩提加雅等）（待核）
\item ★ 【辑1】\textbf{宝鸡——广元} — 宝鸡（陕）—广元（川），蜀道（待核）
\item ○ 【辑1】\textbf{过徐州} — 徐州（待核）
\item ★ 【辑1】\textbf{成都散记} — 成都（待核）
\item ★ 【辑1】\textbf{昆明杂记} — 昆明（待核）
\item ○ 【辑1】\textbf{森林·雨季·山头人} — 西南/滇缅印一带（待核）（待核）
\item ○ 【辑1】\textbf{音尘} — 西南后方（待核）（推测）
\item ○ 【辑1】\textbf{茶馆} — 西南（成渝一带，待核）（待核）
\item ○ 【辑1】\textbf{江上杂记} — 长江（舟行，具体江段待核）（待核）
\item ○ 【辑1】\textbf{断片} — （追忆，非行旅地点）（待核）
\end{itemize}


\subsection{约1944—1947}


\begin{itemize}
\item ○ 【辑2】\textbf{前线景象} — 前线（待核，疑滇西/缅北）（待核）
\item ★ 【辑2】\textbf{他们怎样看中国} — 西南后方（待核）（待核）
\item · 【辑2】\textbf{“中国通”} — 西南后方（待核）（待核）
\item ○ 【辑2】\textbf{人种·职业·人性的大集合} — 西南后方（待核）（待核）
\item ○ 【辑2】\textbf{几个人物} — 西南后方（待核）（待核）
\item ★ 【辑2】\textbf{中国将军怎样应付美国兵} — 西南后方（待核）（待核）
\item ★ 【辑2】\textbf{美国兵与女人} — 西南后方（待核）（待核）
\item · 【辑2】\textbf{“伟大”的S.O.S.} — 西南后方（待核）（待核）
\item · 【辑2】\textbf{咖啡与战斗力} — 西南后方（待核）（待核）
\item ○ 【辑2】\textbf{种种惊异} — 西南后方（待核）（待核）
\item · 【辑2】\textbf{军中文化} — 西南后方（待核）（待核）
\item ★ 【辑2】\textbf{为美国兵活着的人们} — 西南后方（待核）（待核）
\item ★ 【辑2】\textbf{关于“翻译官”} — 西南后方（待核）（待核）
\item · 【辑2】\textbf{汽车团} — 西南后方（待核，疑滇缅公路沿线）（待核）
\end{itemize}


\subsection{1946}


\begin{itemize}
\item ○ 【辑3】\textbf{旅京随笔} — 南京（待核）
\item ○ 【辑3】\textbf{关于“泽存书库”} — 南京（待核）
\item ○ 【辑3】\textbf{访“盔山精舍”} — 南京（待核）
\item ○ 【辑3】\textbf{「美人肝」} — 南京（待核）
\item ★ 【辑3】\textbf{老虎桥边看“知堂”} — 南京·老虎桥监狱（待核）
\item ★ 【辑3】\textbf{鸡鸣寺} — 南京·鸡鸣寺（待核）
\end{itemize}


\subsection{1947}


\begin{itemize}
\item ★ 【辑3】\textbf{金陵杂记} — 南京（多处）（待核）
\end{itemize}


\subsection{约1947}


\begin{itemize}
\item · 【辑2】\textbf{叙言} — —（待核）
\end{itemize}


\subsection{1949}


\begin{itemize}
\item ○ 【辑3】\textbf{司徒雷登夹着皮包走了以后} — 南京（待核）
\item ○ 【辑3】\textbf{社会科学者眼里的南京} — 南京（待核）
\item ○ 【辑3】\textbf{不再做花瓶——记中国农业机械公司南京工厂} — 南京（待核）
\item ★ 【辑3】\textbf{解放后看江南} — 江南／南京（待核）
\end{itemize}


\subsection{1979}


\begin{itemize}
\item ○ 【辑3】\textbf{南唐二陵} — 南京·南唐二陵（待核）
\item ★ 【辑3】\textbf{天王府的西花园} — 南京·天王府（太平天国旧址）（待核）
\item ○ 【辑3】\textbf{扫叶楼} — 南京·扫叶楼（清凉山）（待核）
\item ○ 【辑3】\textbf{梅园（1979版）} — 南京·梅园（待核）
\item ○ 【辑3】\textbf{白鹭洲公园补记} — 南京·白鹭洲（待核）
\item ○ 【辑3】\textbf{石巢园} — 南京·石巢园（待核）
\item ★ 【辑3】\textbf{秦淮拾梦记} — 南京·秦淮（待核）
\item ○ 【辑3】\textbf{莫愁湖（1979版）} — 南京·莫愁湖（待核）
\item ★ 【辑3】\textbf{重过鸡鸣寺} — 南京·鸡鸣寺（待核）
\item ○ 【辑3】\textbf{献花岩} — 南京郊外·献花岩（具体地点待核）（待核）
\item ★ 【辑3】\textbf{王介甫与金陵} — 南京（半山园等王安石相关遗迹）（待核）
\item ○ 【辑3】\textbf{白下书简} — 南京（白下）（待核）
\end{itemize}


\subsection{年份待核}


\begin{itemize}
\item ★ 【辑8】\textbf{余淡心与金陵} — 南京（秦淮）
\item ○ 【辑8】\textbf{关于余淡心} — —（涉南京秦淮等余怀相关地点）
\item ○ 【辑8】\textbf{关于吴梅村} — —（涉江南，吴伟业相关行迹）
\item ○ 【辑8】\textbf{咏怀堂} — —（涉江苏如皋或南京等冒辟疆相关地点，待核）
\item ○ 【辑4】\textbf{杭州杂记} — 杭州
\item ○ 【辑4】\textbf{海滨消夏记} — 海滨（具体地待核，疑青岛/北戴河一带）
\item ★ 【辑5】\textbf{琉璃厂} — 北京·琉璃厂
\item ○ 【辑8】\textbf{绝代的散文家——张宗子} — —（涉绍兴/杭州等张岱行迹）
\item ★ 【辑7】\textbf{西泠访书记} — 杭州·西泠（西湖孤山一带）
\item ○ 【辑4】\textbf{采石·当涂·青山} — 采石矶、当涂、青山（皖南／长江沿线）
\end{itemize}


\medskip\hrule\medskip

\subsection{各辑篇目统计}


总计 59 篇


\begin{itemize}
\item 第1辑 蜀道与战时中国：9 篇
\item 第2辑 美国兵与异国战争经验：16 篇
\item 第3辑 金陵旧梦：24 篇
\item 第4辑 江南行旅：3 篇
\item 第5辑 北平与旧都余韵：1 篇
\item 第7辑 书肆、旧书与纸上行旅：1 篇
\item 第8辑 人物途中：5 篇
\end{itemize}


\section{附录六　待核清单与已知缺口}
\subsection{待核清单（needs\_verification）}
\begin{itemize}
\item 【已更新】《锦帆集》丛刊归属已多源确认为'中华文艺丛刊'；《锦帆集外》已多源确认为'文学丛刊第九集（巴金主编）'；仍待图书馆书目最终确认。
\item 【已更新】《金陵五记》出版社已更正为'金陵书画社'（1982年6月，统一书号10234-004）；旧录'江苏人民出版社'系误，已删；仍待实物核定。《花步集》《晚春的行旅》出版社/年份仍未坐实。
\item 【已更新】《昆明杂记》已更正归属《锦帆集外》（非《锦帆集》）；五部分结构已多源确认；仍待原书核定。
\item 【已更新】《音尘集》（辽宁教育出版社，1996，书趣文丛第三辑）及《印度小夜曲》已录入；年份、出版社待实物核。
\item 《白门秋柳》单篇（1942，初出《锦帆集》ch3）与同名后出选本（豆瓣26801544）须严格区分。白门秋柳是否存在'锦帆集本'与'金陵五记本'异文，待原书比对。
\item 关于美国兵各篇的确切写作时间与发生地点（多笼统标'西南后方/前线'，待据原文定位）。
\item 《重过鸡鸣寺》（E-012）写作年份：《金陵五记》编年列1979，一说1980；待以原书内文日期核定。
\item 西泠访书记、琉璃厂、海滨消夏记的写作年份与原书归属仍待核。
\item 《关于美国兵》15篇个别字词（如'S.O.S.'所指建制）须以原书校。
\item 印度小夜曲八章内容须以《音尘集》原书核定（八章章名已多源确认，内容细节待核）。
\end{itemize}
\subsection{已知缺口（gaps）}
\begin{itemize}
\item 第三辑（金陵）：《金陵五记》五个时间层已基本录入（59篇中24篇属第三辑）；梅园(1979)与莫愁湖(1979)须核定是否与金陵杂记子目重合；「美人肝」内容完全待核。
\item 《锦帆集外》：《昆明杂记》（五部分）已录入（E-018）；其余集外文字（除《昆明杂记》外）完整篇目尚待录入。
\item 第二辑新增：《印度小夜曲》（E-060，《音尘集》1996）已录入为'必收'；八章结构确认，内容细节待原书核定。
\item 第四辑（江南）：仍仅3篇，《花步集》《晚春的行旅》逐篇目录获取受阻（豆瓣/中图网 403），须改走图书馆 OPAC（国图/上图）或《黄裳集·创作卷IV》原书。
\item 第五辑（北平/琉璃厂/旧城/剧场）仅《琉璃厂》1篇，须从《翠墨集》《珠还记幸》《来燕榭文存》补采旧城/剧场/访书诸篇。
\item 第六辑（古迹/废墟，跨地域）独立篇目为零；天王府等怀古篇因金陵地缘划入第三辑；须从《翠墨集》'山川风物'内容（目录待录入）补充跨地域古迹篇。
\item 戏曲题材（剧场/戏曲人物）篇目几乎空白，须查《旧戏新谈》中含地点/人物的篇目，选入第五辑或第八辑。
\item 第七辑（书旅）仅《西泠访书记》1篇；须从《翠墨集》《榆下杂说》《珠还记幸》中摘选带行踪的访书文。
\item 第八辑（人物）现5篇；《珠还记幸》《来燕榭文存》中记人文字尚未录入。
\item 多部原始结集的出版社/年份仍为网络二手，须以图书馆书目/实物核定（金陵五记出版社已更正，其他仍待核）。
\item 尚未系统排查'同篇异文'（如白门秋柳的锦帆集本 vs 金陵五记本；梅园/莫愁湖1947 vs 1979版）。
\end{itemize}
\subsection{不宜收判例（not\_recommended）}
\begin{itemize}
\item \textbf{纯书话（无行踪）：谈"善本"、谈"题跋"、谈"集部"、谈禁书、古书的作伪 等（《榆下说书》）}　属版本学/书话专论，无地点或行旅线索，与'行旅文史'主旨偏离；留待黄裳书话专集，不入本选本。
\item \textbf{纯文学评论：关于张宗子 / 关于吴梅村 / 关于余淡心 等单纯论人论文之作}　论人论文成分重、行旅成分弱；除非与具体地点（如余淡心与金陵）绑定，否则列'不宜收'或'可选'，避免冲淡行旅主线。
\item \textbf{纯时政评论篇（如部分1949年时评）}　时评性强、时效性内容多，须经史地注释方能入；超出'地点触发的文史随笔'范畴者不收。
\end{itemize}

\end{document}

